\documentclass[10pt]{letter}
\usepackage[margin=0.8in]{geometry}
\usepackage{hyperref}
\hypersetup{colorlinks=true, urlcolor=blue}
\setlength{\emergencystretch}{3em}

\longindentation=0pt

\signature{Albert Linda}

\address{Albert Linda \\
Department of Materials Science and Engineering \\
Indian Institute of Technology Kanpur, Kanpur 208016, India \\
\texttt{albert.hzbn@gmail.com}}

\date{\today}

\begin{document}

\begin{letter}{The Editor \\
\textit{Physical Review B} \\
American Physical Society}

\opening{Dear Editor,}

I am pleased to resubmit the following revised manuscript for consideration in
\textit{Physical Review B}:

\noindent\textbf{``Special Optimal Thermal Cell: Correlator-Designed
Supercells for Effective-Harmonic Phonons and Thermodynamics''}

\noindent Albert Linda

We thank the reviewers for identifying several important issues in the original
submission.  The revision substantially narrows the claims and adds direct,
auditable controls.  SOTC is now presented as deterministic configuration
design followed by symmetry-constrained effective second-order force-constant
(FC2) fitting.  It is not presented as a variational anharmonic theory or as a
general replacement for TDEP, SSCHA, or molecular dynamics.

The revised manuscript reports nine metal and rocksalt benchmarks.  Their
archived workflows reach the scalar correlator stopping criterion in 2--5
iterations using 4--36 force calculations, including reference structures;
these counts are reported as measurements rather than as a universal speedup.
A new matched Quantum ESPRESSO comparison for Al, W, and NaCl evaluates 20
canonical-Gaussian sequences and three pre-specified correlator-matched
sequences at identical force-call counts.  At 16 calls the nominal SOTC
dispersion errors differ from the canonical medians by $+3.2\%$, $+8.1\%$, and
$-2.8\%$, respectively.  We therefore conclude only that the nominal design is
comparable in these tests and explicitly report its initialization sensitivity.

The revision also corrects the factor-of-three mean-square-displacement
convention, removes the underconverged quasi-harmonic calculation and all
dependent claims, replaces unsupported rank statements with auditable
force-call and saved-matrix diagnostics, and defines reference-force
subtraction explicitly.  New calculations provide non-analytic LO--TO
corrections for NaCl, MgO, and PbTe; tabulated 80~K NaCl neutron data are shown
directly; and the Au instability is quantified and excluded from the stable
dispersion evidence.  The response letter gives a point-by-point account of
these changes.

The revised source, response letter, validation scripts, compact
machine-readable evidence, and generated figures are available in the
\href{https://github.com/albert-hzbn/SOTC}{public SOTC repository}.

The manuscript is a good fit for \textit{Physical Review B}: it presents a
first-principles method for condensed-matter thermodynamics, supported by
theoretical derivations, systematic validation, and open reproducibility.

This manuscript has not been submitted elsewhere and is not under
consideration at any other journal. I have no conflicts of interest to
declare.

Thank you for your time and consideration.

\closing{Sincerely,}

\end{letter}
\end{document}
